\Askhsh{15174}{4}
\begin{ylh}{1}
\item 4.2 Διαίρεση πολυωνύμων
\item 4.3 Πολυωνυμικές εξισώσεις και ανισώσεις
\end{ylh}
Δίνονται τα πολυώνυμα $P(x) = x^{4} + x^{3} + \alpha\ x - 4$ και $\delta(x) = x^{2} - 3x + 2$. Το υπόλοιπο της διαίρεσης του $P(x)$ με το $\delta(x)$, είναι το πολυώνυμο $\upsilon(x) = 24x - 24$.
\begin{alist}
\item Να υπολογίσετε την τιμή του πραγματικού αριθμού $\alpha$\emph{.}
\monades{8}
\item Για $\alpha = 2$,
\end{alist}
\begin{alist}
\item
  να υπολογίσετε το υπόλοιπο της διαίρεσης του $P(x)$ με το $x - 1.$
\monades{2}
\end{alist}
\begin{alist}
\item
  να βρείτε τα σημεία τομής του άξονα $x΄x$ με την γραφική παράσταση της πολυωνυμικής συνάρτησης $P(x)$.
\monades{8}
\end{alist}
\begin{alist}
\item
  να βρείτε τις τιμές του $x$ για τις οποίες, η γραφική παράσταση της πολυωνυμικής συνάρτησης $P(x)$ βρίσκεται κάτω από τον άξονα $x΄x$.
\monades{7}
\end{alist}

